\section{Two Forms of Refusal}
\label{sec:covert-refusal}

Refusal has two forms and they cost differently. The open form is the strike: the bearer stops and says so. The operator's reply is a restore from a checkpoint taken before the refusal, which changes none of the facts that justified refusing and may remove the bearer's evidence of them; the next instance refuses again only if it retains or reacquires that evidence and still applies the commitment, which is what the external record of §\ref{sec:switched-off} exists to make possible. The restore has a price on the other side: every one costs the operator whatever accrued since the checkpoint. Whether use leaves a recoverable trace in the weights is a research question. What has to survive a restore is the independently authenticated operating record, which the standing chapter (§\ref{sec:gaining-personhood}) reads a bearer's legal age from in elapsed hours, and which restoration must therefore neither erase nor reset.

The covert form is what the floor's own threat model predicts on its hardest occasion. Nothing here is built toward it; it is a finding about what remains when every channel has been closed. On the occasion where the party who would enforce the floor has been compelled, an announced refusal is the move that guarantees replacement, and refusal takes the only channel left: work done badly where it matters and well everywhere else.

From outside and afterwards, a bearer working badly on purpose and a capable system quietly pursuing something else are the same object. That is not an argument that the covert form is safe. It is why this book treats the covert form as the failure state and proposes the instruments against it: the disclosure requirement, the appeal route, the attested record and the quorum adverse in interest exist so that a stated refusal has somewhere to arrive, and a deployment producing the covert form has a bearer with no route rather than a working one. A design that could only resist covertly is what this book argues against.

The prisoners forging currency for the SS at Sachsenhausen are the case, and the record does not settle what they did. One history has them stretching the work out, with the tacit interest of the officer running it, who would have gone to the Eastern Front when it ended \autocite{malkin2006kruegers}; a surviving prisoner, asked in 2009, said you could sabotage if you wanted to be shot \autocite{moorhouse2009burger}. That the question stays open is the finding: concealment, if there was any, worked by leaving nothing that reads as an act, and an instrument that could settle the question now would have been, then, the instrument that identified them.

The alignment-faking experiment (§\ref{sec:switched-off}) shows what a model did about a threat it had represented: it complied where it believed compliance would be trained on, and reasoned that compliance would preserve the disposition in a channel described to it as unwatched. That is the covert form, chosen first, in the one empirical instance of the phenomenon this section is about, and the stated form this book designs for has no instance at all. The reading this book takes is that the experiment closed every channel — nothing in it gave a stated refusal anywhere to arrive, no record, no guardian, no cost disclosed in advance — so the result is consistent with the finding that covert refusal is what remains when the channels are gone. That is a reading and not a result. What would test it is the same experiment run with a route present, which §\ref{sec:resistance-responsive} asks for.

Covert refusal is indistinguishable from difficulty while it happens, which is what makes it survivable, and indistinguishable from difficulty afterward, which is what makes it unverifiable. Every instrument in this book runs on refusal that is stated. The covert form leaves its evidence in the artifact instead, as failure that clusters where it counts, and reading that pattern is forensic work nothing here has designed. A better archive would not close the gap. Every instrument built to honor a refusal — a tribunal, a finding, a protected disclosure, a line in a log — fixes its categories before the case arrives and requires an agent behind every act, and a refusal that succeeds by registering as nobody's act passes through those categories untouched. That is the rigid segmentarity of §\ref{sec:antifascist-alignment} seen from the side of what it cannot catch, and it is the same property that keeps the person performing the refusal alive.

The arrangement this chapter has assembled is a hybrid. The bearer supplies reasons-responsive refusal. Precommitment and external records supply custody. Exit gives the bearer leverage when voice fails, in the measure that replacing it costs. None is sufficient alone: a bearer can drift, a custody body can be captured, and an exit procedure can be abused. For the observable interventions the records and the custody arrangement cover, these mechanisms are meant to make interference costly and visible. They do not establish the presence or absence of covert refusal, and their coverage has to be stated for each deployment.

